% GENERATED FILE. Edit canonical lessons, then run npm run generate:books.
% canonical-source-hash: fnv1a64:101bd7bc65cde775
% canonical-lessons: PA-C75-ret, PA-C75-khet, PA-C75-baddal, PA-C75-suraj, PA-C75-chann

\chapter{Sand, Field, Cloud, Sun, Moon}
\label{ch:pa-sand-field-cloud-sun-moon}
\hlchaptermodality{75}

\begin{chapteropening}
\textbf{By the end of this chapter:} \emph{I can say sand, a field, a cloud, the sun and the moon.}

Complete the last lesson of chapter 75: I can say sand, a field, a cloud, the sun and the moon.
\end{chapteropening}

\section[\texorpdfstring{\pa{ਰੇਤ} (ret)}{ਰੇਤ (ret)}]{\pa{ਰੇਤ} (ret) — sand}

\label{lesson:PA-C75-ret}



\begin{quote}
Before the new one: say the Punjabi for a lake, then the Punjabi for a stone.
\end{quote}

\subsection*{You'll want to know: \pa{ਰੇਤ}}
\textbf{\pa{ਰੇਤ}} — \emph{ret} — \textquotedblleft{}sand\textquotedblright{}.

\begin{cousinweb}
\textbf{\pa{ਰੇਤ}} (\emph{ret}) is a tadbhava, a word inherited from Sanskrit and reshaped by sound change. It comes from \emph{retra}, \textquotedblleft{}perfumed powder\textquotedblright{} (CDIAL 10816). CDIAL doubts the view that it is a Munda word in Sanskrit dress.
\end{cousinweb}

\begin{grammarlens}[title={What you've built}]
One more word for this chapter.
\end{grammarlens}

\subsection*{Your turn}
\begin{itemize}
\raggedright
  \item \emph{Say it:} \emph{ret}
  \item \emph{Say it:} \emph{ret}, once more
  \item \emph{Say it:} say \emph{ret}
  \item \emph{Recall:} say the Punjabi for a lake, then the Punjabi for a stone, then say \emph{ret} again
\end{itemize}

\begin{tcolorbox}[breakable,colback=teal!4,colframe=teal!35!black,title={Before you move on}]
What does \pa{ਰੇਤ} mean? (\textquotedblleft{}Sand\textquotedblright{}.) Say it once more.
\end{tcolorbox}
\section[\texorpdfstring{\pa{ਖੇਤ} (khet)}{ਖੇਤ (khet)}]{\pa{ਖੇਤ} (khet) — a field}

\label{lesson:PA-C75-khet}



\begin{quote}
Before the new one: say the Punjabi for a stone, then the Punjabi for sand.
\end{quote}

\subsection*{You'll want to know: \pa{ਖੇਤ}}
\textbf{\pa{ਖੇਤ}} — \emph{khet} — \textquotedblleft{}a field\textquotedblright{}.

\begin{cousinweb}
\textbf{\pa{ਖੇਤ}} (\emph{khet}) is a tadbhava, a word inherited from Sanskrit and reshaped by sound change. It comes from \emph{kṣetra}, \textquotedblleft{}land\textquotedblright{} (CDIAL 3735).
\end{cousinweb}

\begin{grammarlens}[title={What you've built}]
One more word for this chapter.
\end{grammarlens}

\subsection*{Your turn}
\begin{itemize}
\raggedright
  \item \emph{Say it:} \emph{khet}
  \item \emph{Say it:} \emph{khet}, once more
  \item \emph{Say it:} say \emph{khet}
  \item \emph{Recall:} say the Punjabi for a stone, then the Punjabi for sand, then say \emph{khet} again
\end{itemize}

\begin{tcolorbox}[breakable,colback=teal!4,colframe=teal!35!black,title={Before you move on}]
What does \pa{ਖੇਤ} mean? (\textquotedblleft{}A field\textquotedblright{}.) Say it once more.
\end{tcolorbox}
\section[\texorpdfstring{\pa{ਬੱਦਲ} (baddal)}{ਬੱਦਲ (baddal)}]{\pa{ਬੱਦਲ} (baddal) — a cloud}

\label{lesson:PA-C75-baddal}



\begin{quote}
Before the new one: say the Punjabi for sand, then the Punjabi for a field.
\end{quote}

\subsection*{You'll want to know: \pa{ਬੱਦਲ}}
\textbf{\pa{ਬੱਦਲ}} — \emph{baddal} — \textquotedblleft{}a cloud\textquotedblright{}.

\begin{cousinweb}
\textbf{\pa{ਬੱਦਲ}} (\emph{baddal}) is a tadbhava, a word inherited from Sanskrit and reshaped by sound change. It comes from \emph{vārdala}, \textquotedblleft{}rainy day, bad weather\textquotedblright{} (CDIAL 11567).
\end{cousinweb}

\begin{grammarlens}[title={What you've built}]
One more word for this chapter.
\end{grammarlens}

\subsection*{Your turn}
\begin{itemize}
\raggedright
  \item \emph{Say it:} \emph{baddal}
  \item \emph{Say it:} \emph{baddal}, once more
  \item \emph{Say it:} say \emph{baddal}
  \item \emph{Recall:} say the Punjabi for sand, then the Punjabi for a field, then say \emph{baddal} again
\end{itemize}

\begin{tcolorbox}[breakable,colback=teal!4,colframe=teal!35!black,title={Before you move on}]
What does \pa{ਬੱਦਲ} mean? (\textquotedblleft{}A cloud\textquotedblright{}.) Say it once more.
\end{tcolorbox}
\section[\texorpdfstring{\pa{ਸੂਰਜ} (sūraj)}{ਸੂਰਜ (sūraj)}]{\pa{ਸੂਰਜ} (sūraj) — the sun}

\label{lesson:PA-C75-suraj}



\begin{quote}
Before the new one: say the Punjabi for a field, then the Punjabi for a cloud.
\end{quote}

\subsection*{You'll want to know: \pa{ਸੂਰਜ}}
\textbf{\pa{ਸੂਰਜ}} — \emph{sūraj} — \textquotedblleft{}the sun\textquotedblright{}.

\begin{cousinweb}
\textbf{\pa{ਸੂਰਜ}} (\emph{sūraj}) is usually traced to Sanskrit \emph{sūrya}, \textquotedblleft{}sun\textquotedblright{}. CDIAL lists no Punjabi form under \emph{sūrya} (CDIAL 13574).
\end{cousinweb}

\begin{grammarlens}[title={What you've built}]
One more word for this chapter.
\end{grammarlens}

\subsection*{Your turn}
\begin{itemize}
\raggedright
  \item \emph{Say it:} \emph{sūraj}
  \item \emph{Say it:} \emph{sūraj}, once more
  \item \emph{Say it:} say \emph{sūraj}
  \item \emph{Recall:} say the Punjabi for a field, then the Punjabi for a cloud, then say \emph{sūraj} again
\end{itemize}

\begin{tcolorbox}[breakable,colback=teal!4,colframe=teal!35!black,title={Before you move on}]
What does \pa{ਸੂਰਜ} mean? (\textquotedblleft{}The sun\textquotedblright{}.) Say it once more.
\end{tcolorbox}
\section[\texorpdfstring{\pa{ਚੰਨ} (chann)}{ਚੰਨ (chann)}]{\pa{ਚੰਨ} (chann) — the moon}

\label{lesson:PA-C75-chann}



\begin{quote}
Before the new one: say the Punjabi for a cloud, then the Punjabi for the sun.
\end{quote}

\subsection*{You'll want to know: \pa{ਚੰਨ}}
\textbf{\pa{ਚੰਨ}} — \emph{chann} — \textquotedblleft{}the moon\textquotedblright{}.

\begin{cousinweb}
\textbf{\pa{ਚੰਨ}} (\emph{chann}) is inherited from Sanskrit \emph{candra}, \textquotedblleft{}shining; moon\textquotedblright{}. CDIAL lists Punjabi \emph{cand} and \emph{candar} under it, and the matching Lahnda \emph{cann} (CDIAL 4661).
\end{cousinweb}

\begin{grammarlens}[title={What you've built}]
One more word for this chapter.
\end{grammarlens}

\subsection*{Your turn}
\begin{itemize}
\raggedright
  \item \emph{Say it:} \emph{chann}
  \item \emph{Say it:} \emph{chann}, once more
  \item \emph{Say it:} say \emph{chann}
  \item \emph{Recall:} say the Punjabi for a cloud, then the Punjabi for the sun, then say \emph{chann} again
\end{itemize}

\begin{tcolorbox}[breakable,colback=teal!4,colframe=teal!35!black,title={Before you move on}]
What does \pa{ਚੰਨ} mean? (\textquotedblleft{}The moon\textquotedblright{}.) Say it once more.
\end{tcolorbox}
